% Part: intuitionistic-logic
% Chapter: tableaux
% Section: rules

\documentclass[../../../include/open-logic-section]{subfiles}

\begin{document}

\olfileid[hi]{int}{tab}{rul}

\olsection{अंतःप्रज्ञावादी तर्कशास्त्र के नियम}

संयोजकों $\land$ और $\lor$ के नियम सामान्य प्रतिज्ञप्तिक चिह्नित
!!{tableau} जैसे ही हैं, केवल उपसर्ग जोड़ दिए गए हैं। प्रत्येक मामले में,
किसी चिह्नित !!{formula} $\sFmla{S}{!A}[\sigma]$ पर लगाया गया नियम ऐसे नए
!!{formula} उत्पन्न करता है जिन पर भी~$\sigma$ उपसर्ग लगा है। यह सहजतः स्पष्ट
होना चाहिए: उदाहरणार्थ यदि $!A \land !B$,~$\sigma$ (द्वारा नामित जगत) पर
सत्य है, तो $!A$ और $!B$,~$\sigma$ पर सत्य हैं (किसी अन्य जगत पर नहीं)।
$\land$ और $\lor$ के नियम \olref{tab:prop-rules} में संकलित हैं।

\begin{table}
  \[\def\arraystretch{3}\begin{array}{|c|c|}
    \hline
    \AxiomC{\sFmla{\True}{!A \land !B}[\sigma]}
    \RightLabel{\TRule{\True}{\land}}
    \UnaryInfC{\sFmla{\True}{!A}[\sigma]}
    \noLine
    \UnaryInfC{\sFmla{\True}{!B}[\sigma]}
    \DisplayProof
    &
    \AxiomC{\sFmla{\False}{!A \land !B}[\sigma]}
    \RightLabel{\TRule{\False}{\land}}
    \UnaryInfC{$\sFmla{\False}{!A}[\sigma] \quad \mid \quad
      \sFmla{\False}{!B}[\sigma]$}
    \DisplayProof
    \\[2ex]
    \hline
    \AxiomC{\sFmla{\True}{!A \lor !B}[\sigma]}
    \RightLabel{\TRule{\True}{\lor}}
    \UnaryInfC{$\sFmla{\True}{!A}[\sigma] \quad \mid \quad
      \sFmla{\True}{!B}[\sigma]$}
    \DisplayProof
    &
    \AxiomC{\sFmla{\False}{!A \lor !B}[\sigma]}
    \RightLabel{\TRule{\False}{\lor}}
    \UnaryInfC{\sFmla{\False}{!A}[\sigma]}
    \noLine
    \UnaryInfC{\sFmla{\False}{!B}[\sigma]}
    \DisplayProof
    \\[2ex]
    \hline
  \end{array}\]
  \caption{$\land$ और $\lor$ के उपसर्गयुक्त !!{tableau} नियम}
  \ollabel{tab:prop-rules}
\end{table}

संवरण-शर्त सामान्य !!{tableau} की शर्त जैसी है, यद्यपि हम यह अपेक्षा करते
हैं कि केवल !!{formula} ही नहीं, उपसर्ग भी मेल खाएँ। वास्तव में हम कुछ अधिक
उदार हो सकते हैं: चूँकि अंतःप्रज्ञावादी प्रतिरूपों में !!{formula} एक बार
सत्य होने पर सत्य रहते हैं, इसलिए यह असंभव है कि $!A$,~$\sigma$ पर सत्य
पर किसी अभिगम्य उपसर्ग~$\sigma.{*}$ पर असत्य हो। अतः कोई शाखा संवृत है यदि
उसमें दोनों
\[
\sFmla{\True}{!A}[\sigma] \quad\text{और}\quad \sFmla{\False}{!A}[\sigma.{*}]
\]
किसी उपसर्ग $\sigma$ और !!{formula}~$!A$ के लिए हों। ध्यान दें कि यदि चिह्न
उलट दिए जाएँ, अर्थात उसमें
\[
\sFmla{\False}{!A}[\sigma] \quad\text{और}\quad \sFmla{\True}{!A}[\sigma.{*}]
\]
हों, तो शाखा केवल तभी संवृत है जब $*$ रिक्त अनुक्रम हो।

इसके अतिरिक्त, यदि किसी शाखा में~$\sFmla{\True}{\bot}[\sigma]$ हो, तो वह
संवृत है।

मान्यताएँ स्थापित करने के नियम भी सामान्य !!{tableau} जैसे हैं, सिवाय इसके
कि मान्यताओं के लिए हम सदैव उपसर्ग~$1$ का उपयोग करते हैं। (कौन-सा उपसर्ग
लेते हैं इससे अंतर नहीं पड़ता, बशर्ते सभी मान्यताओं के लिए वही हो।) अतः,
उदाहरणार्थ, हम कहते हैं कि
\[
!B_1, \dots, !B_n \Proves !A
\]
तभी जब निम्न मान्यताओं के लिए कोई संवृत टैब्लो हो:
\[
\sFmla{\True}{!B_1}[1], \dots, \sFmla{\True}{!B_n}[1],
\sFmla{\False}{!A}[1].
\]

सशर्त~$\lif$ के नियम शास्त्रीय और मोडल मामलों से भिन्न हैं।
$\TRule{\lif}{\True}$ नियम, $\sFmla{\True}{!A \lif !B}[\sigma]$ वाली शाखा
को दो अलग शाखाओं पर $\sFmla{\True}{!A}[\sigma.{*}]$ और
$\sFmla{\False}{!B}[\sigma.{*}]$ जोड़कर विस्तृत करता है। इसे केवल ऐसे
उपसर्ग~$\sigma.{*}$ के लिए लगाया जा सकता है जो उस शाखा पर \emph{पहले से}
आता हो। ऐसे उपसर्ग को शाखा पर ``प्रयुक्त'' कहें। (चूँकि $\sigma.{*}$ में
स्वयं $\sigma$ भी आता है, नियम को अलग शाखाओं पर उपसर्गयुक्त चिह्नित सूत्र
$\sFmla{\True}{!A}[\sigma]$ और $\sFmla{\False}{!B}[\sigma]$ जोड़कर सदैव
लगाया जा सकता है।)

$\TRule{\lif}{\False}$ नियम, $\sFmla{\False}{!A \lif !B}[\sigma]$ वाली
शाखा को उसी शाखा पर $\sFmla{\True}{!A}[\sigma.n]$ और
$\sFmla{\False}{!B}[\sigma.n]$ दोनों जोड़कर विस्तृत करता है, जहाँ
$\sigma.n$ उस शाखा के लिए नया उपसर्ग है।

~$\lnot$ के नियम इसी के अनुरूप परिभाषित हैं ($\lnot !A$ की
$!A \lif \lfalse$ वाली परिभाषा का उपयोग करके)।

नियम \olref{tab:rules-lif-lnot} में दिए गए हैं।

\begin{table}
  \[\def\arraystretch{3}\begin{array}{|c|c|}
    \hline
    \AxiomC{\sFmla{\True}{\lnot !A}[\sigma]}
    \RightLabel{\TRule{\True}{\lnot}}
    \UnaryInfC{$\sFmla{\False}{!A}[\sigma.{*}]$}
    \DisplayProof
    &
    \AxiomC{\sFmla{\False}{\lnot !A}[\sigma]}
    \RightLabel{\TRule{\False}{\lnot}}
    \UnaryInfC{\sFmla{\True}{!A}[\sigma.n]}
    \DisplayProof
    \\[1ex]
    \text{$\sigma.{*}$ प्रयुक्त है} & \text{$\sigma.n$ नया है}\\
    \hline
    \AxiomC{\sFmla{\True}{!A \lif !B}[\sigma]}
    \RightLabel{\TRule{\True}{\lif}}
    \UnaryInfC{$\sFmla{\False}{!A}[\sigma.{*}] \quad \mid \quad
      \sFmla{\True}{!B}[\sigma.{*}]$}
    \DisplayProof
    &
    \AxiomC{\sFmla{\False}{!A \lif !B}[\sigma]}
    \RightLabel{\TRule{\False}{\lif}}
    \UnaryInfC{\sFmla{\True}{!A}[\sigma.n]}
    \noLine
    \UnaryInfC{\sFmla{\False}{!B}[\sigma.n]}
    \DisplayProof
    \\[1ex]
    \text{$\sigma.{*}$ प्रयुक्त है} & \text{$\sigma.n$ नया है}\\
    \hline
  \end{array}\]
  \caption{$\lnot$ और $\lif$ के उपसर्गयुक्त !!{tableau} नियम}
  \ollabel{tab:rules-lif-lnot}
\end{table}

\end{document}
